\section{Discussion} \label{sec:discussion}
In this paper, we have comprehensively demonstrated how \sysname improves the performance of LLM inference across several models and hardware configurations. However, there are multiple challenges that require further investigation.

First, we focus only on an efficient scheduling mechanism in \sysname to improve the throughput of LLM inference. However, real-world deployments need to optimize an inference serving infrastructure simultaneously along multiple dimensions e.g., latency, queuing delays, fairness, etc. Meeting these goals with \sysname requires revisiting scheduling policies. Second, although we show what is an appropriate chunk size for a given P:D ratio, we leave it to future work to explore how to pick an optimal chunk size as it depends on several factors like the hardware, model characteristics, sequence length, and the composition of prefill-decode tokens, especially in scenarios where the P:D ratio may not be known ahead of time. Third, we make a simplistic assumption in this paper that each request in a batch has the same number of prefill and decode tokens (except the simulation experiments) whereas, in the real world, the sequence lengths can vary significantly across different LLM inference requests. Finally, we focused on sequence lengths of up to 3K, and P:D ratio in the range of 1-200. We believe that these are representative of many real-world deployments. However, there has also been an increased interest in supporting very long sequences (e.g., 10s-100s of thousands~\cite{largecontextlength}). Such large sequence lengths may pose new challenges as the cost of attention grows quadratically with the number of tokens. We are actively investigating these challenges. 




